\section{Introduction}
\label{sec:intro}

\looseness -1
Diffusion models are the dominant generative model class for high-dimensional data such as images, video, audio, and molecules~\citep{ho2020denoising, song2021scoresde, rombach2022highresolution, karras2022elucidating, leeGenMolDrugDiscovery2025}. Sampling corresponds to numerically integrating the time-reversed stochastic differential equation associated with a forward noising process; the only learnable component is a score network $s_\theta(x_t, t) \approx \nabla \log p_t(x_t)$ trained to match the noisy marginals of the data distribution. Sampling can equivalently be expressed as a deterministic ordinary differential equation with the score as a velocity field, an approach taken by flow-matching methods~\citep{lipman2023flow}.

\looseness -1
Latent diffusion models~\citep{rombach2022highresolution, esser2024sd3, podell2023sdxl} run forward and backward diffusion not in pixel space but in a learned latent space of much smaller dimensionality $\dlat$, typically obtained from a VAE or similar encoder. By placing diffusion in a representation closer to the data's true intrinsic dimensionality $\dint$~\citep{pope2021, brown2022manifold, stanczuk2024manifold, kamkari2024geometric}, latent diffusion achieves comparable or better sample quality at a fraction of the compute of pixel-space diffusion.

\looseness -1
Memorization is not just an empirical accident of overtrained diffusion models. For a finite training set, the exact score of the noised empirical distribution has a closed form: it is built from kernels centered on the training examples, and following this score back to low noise can return the training samples themselves~\citep{scarvelis2023closedform}. Equivalently, the denoising-score-matching objective has a closed-form optimum that reproduces training data rather than drawing genuinely new samples~\citep{gu2023memorization}. Practical score networks usually generalize because finite capacity, approximation error, optimization dynamics, and data geometry smooth this empirical score; recent work studies this gap through dataset-size thresholds, geometry-adaptive representations, and dynamical transitions between population learning and sample collapse~\citep{kadkhodaie2023learning,biroli2024dynamical,bonnaire2025}.

\looseness -1
This makes memorization both a privacy/copyright problem and a basic question about what diffusion models learn. In the broader generative-model literature, memorization has been studied through membership inference, reconstruction attacks, and data-copying tests that ask whether generated samples are unusually close to training examples~\citep{carlini2019secret,hayes2017logan,hilprecht2019reconstruction,meehan2020datacopying}. For diffusion models specifically, nearest-neighbor replication tests and extraction attacks show that real image generators can reproduce training images, while effective-memorization studies tie this risk to training-set size, model capacity, conditioning, and optimization choices~\citep{somepalli2023diffusion,carlini2023extracting,gu2023memorization}. What remains unclear is how the transition depends on the dimension of the latent space itself, especially when diffusion is trained in VAE latents but memorization is measured after decoding back to pixel space. We ask: \emph{how does the gap between $\dlat$ of the model and $\dint$ of the data shape the transition from generalization to memorization, and is there a way a practitioner can predict this behavior?}

\looseness -1
Existing theory provides a starting point. \citet{bonnaire2025} characterize diffusion training as a competition between a generalization time $\taugen$, at which the score network learns the population distribution, and a memorization time $\taumem$, at which it begins fitting individual training points. Their Theorem 3.1 is stated for an arbitrary input spectral measure, but their experiments and their two-bulk reading take the isotropic case, which collapses the distinction between the $\dint$-dimensional signal subspace and the $\dlat-\dint$ null directions; \citet{george2025dsm} treat data on a $\dint$-dimensional subspace with $\signoise=0$. We evaluate their spectral theory at the two-atom covariance of the latent-diffusion regime $\dlat>\dint$ with strictly positive null variance, identify the noise-dimension block that opens between the signal and sample modes, and show that what delays memorization is the growth with $\dlat$ of the ratio between the slowest signal mode and the fastest sample-specific mode.

\textbf{Contributions.}\par\vspace{-0.75em}
\begin{itemize}
  \setlength{\itemsep}{1pt}
  \setlength{\topsep}{0pt}
  \setlength{\partopsep}{0pt}
  \setlength{\parsep}{0pt}
  \item We evaluate Bonnaire et al.'s random-feature spectral theory at the two-atom latent covariance and obtain a four-block structure of the feature correlation matrix with block boundaries at indices $\dint$, $\dlat$ and $\nsamp$ (signal, noise-dimension, sample-specific, rank-null), with block edges that scale as $a_1(q)^2\alpha_t^2/\dlat$, $a_1(q)^2\beta_t^2/\dlat$ and $a_*(q)^2/\nsamp$ and are independent of the feature width $\pwidth$.
  \item We show the buffer is a ratio of spectral edges, not a mode count: $\taugen$ (set by the smallest signal eigenvalue) grows about $11\times$ in gradient-flow time over $\dlat=5\to200$, but the onset of memorization (set by the largest sample-specific eigenvalue) grows about $220\times$, so their ratio grows monotonically by $20\times$. The driver is the fall of the pre-activation variance $q=\operatorname{tr}M_t/\dlat$, which de-saturates $\tanh$ and shrinks the sample-specific feature mass $a_*(q)^2$ by $17\times$.
  \item We test the mechanism in trainable MLP score networks on synthetic data with five-seed confidence intervals, holding hidden width fixed so the $\dlat$ sweep is not a capacity sweep: the memorized fraction falls from $0.30$ to $\approx 0$ over $\dlat=5\to40$, and an exactly solvable linear score model predicts the accompanying score-error curves without fitted parameters.
  \item We test the same mechanism on CelebA and CIFAR-10 latent diffusion. Memorization, measured by pixel-space nearest neighbors after decoding, decreases with latent width once the VAE is expressive enough to preserve sample identity and the sweep enters the excess-latent regime. FID follows a separate quality tradeoff: it improves as the VAE bottleneck is relieved, then worsens when the latent score problem becomes too wide.
  \item We introduce a way to estimate memorization at a given training step from the trained VAE and the chosen training subset: for each latent width, we encode the same $n$ images with the frozen VAE, compute the covariance spectrum of the latent codes, and convert that spectrum into a mode-absorption clock that outputs an estimated memorization percentage over training time. The theory shows which term of the latent spectrum the clock must carry (the noise-averaged sample-specific eigenvalue $g(q_\mu)/\nsamp$, which the covariance spectrum alone does not contain) and reduces the fit to a single scalar.
\end{itemize}
